% !TEX root = ../../cookbook.tex
\newpage
\recipe{Moussaka}{}{}{}

\begin{multicols}{2}
	\subsection*{Ingredients}
	
	\subsubsection*{For the Meat Sauce}
	\begin{ingredients}
		\item 1 kg (2.2 lb) lean ground meat (\sfrac{1}{2} beef and \sfrac{1}{2} pork)
		\item 100 g (3.5 oz) unseasoned tomato sauce (\textit{passata})
		\item 25 g (1.5 tbsp) tomato paste
		\item 1 onion
		\item 1 celery stalk
		\item 2 garlic cloves
		\item \sfrac{1}{2} tsp sugar
		\item 1 glass (175 ml / \sfrac{3}{4} cup) red wine
		\item 2 bay leaves
		\item 1 cinnamon stick
		\item Star anise, to taste
		\item Parsley, to taste
		\item Nutmeg, to taste
		\item Olive oil, as needed
		\item Salt and black pepper, to taste
	\end{ingredients}
	
	\subsubsection*{For the Layers}
	\begin{ingredients}
		\item 1 kg (2.2 lb) eggplants
		\item If frying:
		\begin{ingredients}
			\item Flour, as needed
			\item Frying oil
		\end{ingredients}
		\item 3 potatoes (optional)
	\end{ingredients}
	
	\subsubsection*{For the Yogurt Béchamel}
	\begin{ingredients}
		\item 500 g (2 cups) milk
		\item 250 g (1 cup) Greek yogurt
		\item 50 g (\sfrac{1}{3} cup) flour
		\item 50 g (3.5 tbsp) butter
		\item Nutmeg, to taste
		\item Black pepper, to taste
	\end{ingredients}
	
	\subsubsection*{To Assemble}
	\begin{ingredients}
		\item Parmesan cheese, to taste
		\item Breadcrumbs, to taste
		\item Dried oregano, to taste
		\item Olive oil, as needed
	\end{ingredients}
	
\end{multicols}

% --- Preparation ---
\textbf{Start with the meat sauce:} \\
Roughly chop the \textbf{celery} and \textbf{onion}.
Sauté them with the \textbf{garlic} until softened, then add the \textbf{sugar}.
Add the \textbf{ground meat} and brown it over high heat, breaking it up with a wooden spoon.
After about 5 minutes, deglaze with the \textbf{red wine} and season with the \textbf{bay leaves}, \textbf{star anise}, and \textbf{cinnamon}.
Stir in the \textbf{tomato passata} and \textbf{tomato paste}, mix well, and let simmer over medium-high heat, then add the chopped \textbf{parsley}.
Season with \textbf{nutmeg}, \textbf{salt}, and \textbf{black pepper}.
Reduce the heat and cook uncovered for approximately 1.5 hours.

\textbf{Prepare the eggplants and potatoes:} \\
Slice the \textbf{eggplants} lengthwise into slices about 1 cm (\sfrac{2}{4} in) thick, and the \textbf{potatoes} into slices about \sfrac{1}{2} cm (\sfrac{1}{4} in) thick.
Traditionally, the vegetables are dredged in flour and fried, but alternatively they may be brushed with olive oil and baked for a simpler, lighter version\footnote{In this case, sprinkle the slices with salt on both sides and let them release some of their moisture for about \sfrac{1}{2} hour, then pat them thoroughly dry before assembling the moussaka.}.
Set them aside.

\textbf{Prepare the yogurt béchamel:} \\
Heat the \textbf{milk} separately.
In a saucepan, melt the \textbf{butter} over medium-low heat and stir in the \textbf{flour} to make a golden roux\footnote{A \textit{roux} is the classic French mixture of equal parts butter and flour cooked together and used as the base for many sauces.}.
Gradually add the hot milk while whisking constantly to eliminate any lumps, then season with \textbf{salt}, \textbf{black pepper}, and \textbf{nutmeg}.
Once you have a thick béchamel, remove it from the heat and let it cool for 5 minutes before stirring in the \textbf{Greek yogurt}, then adjust the seasoning with salt, pepper, and nutmeg.

\textbf{Assemble the moussaka:} \\
Grease a deep baking dish (base approximately 26$\times$20 cm (10$\times$8 in)) with \textbf{olive oil}, spread a thin layer of béchamel over the bottom, arrange all the potatoes in a single layer, add a layer of eggplants, and sprinkle with \textbf{Parmesan cheese}, \textbf{breadcrumbs}, and \textbf{dried oregano}.
Cover with all of the meat sauce.
Arrange another layer of eggplants, then sprinkle again with Parmesan cheese, breadcrumbs, and dried oregano.
Cover everything with the remaining béchamel, then finish with another sprinkling of breadcrumbs and Parmesan cheese.
Bake in a preheated oven at 180°C (350°F) for approximately 45 minutes.

% --- Description ---
\begin{recipeblock}
	
\end{recipeblock}